\documentclass[../../../include/open-logic-section]{subfiles}

\begin{document}

\olfileid{sth}{z}{sep}
\olsection{Pamisahan}

Kita miwiti kanthi prinsip kanggo ngganti Komprehensi Naif:

\begin{axiom}[Skema Pamisahan] Kanggo saben formula $\phi(x)$, iki
aksioma: kanggo sembarang $A$, himpunan $\Setabs{x \in A}{\phi(x)}$ ana.
\end{axiom}

Gatekna yen iki dudu siji aksioma. Iki \emph{skema} aksioma.
Ana aksioma Pamisahan kang \emph{tanpa wates cacahé}; siji
kanggo saben formula $\phi(x)$. Skema iki uga bisa
(lan lumrahe pancen) ditulis kaya mangkene:

\begin{defish}
Kanggo sembarang formula $\phi(x)$ kang ora ngemot ``$S$'',
iki sawijining aksioma:
\[
	\forall A \exists S \forall x(x \in S \liff (\phi(x) \land x \in A)).
\]
\end{defish}

Miturut konvensi kang kasebut ing wiwitan
\olref[sth][][]{part}, formula-formula~$\phi$ ing aksioma
Pamisahan bisa ngemot parameter.\footnote{Kanggo nerangake
tegese, delengen rembugan pas sawise
\olref[sfr][infinite][induction]{natinductionschema}.}

Pamisahan langsung didhukung dening konsepsi
kumulatif-iteratif tumrap himpunan kang wis kita aturake.
Kanggo ndeleng sebabe, upamane $A$ iku himpunan.
Mula $A$ kawangun ing sawenehing tataran~$S$
(miturut \stageshier). Amarga $A$ kawangun ing tataran~$S$,
kabeh anggotane $A$ wis kawangun sadurunge tataran $S$
(miturut \stagesacc). Saiki nimbang mligine kabeh himpunan
kang dadi anggota $A$ lan uga nyukupi $\phi$; cetha yen
himpunan-himpunan iki kabeh wis kawangun sadurunge
tataran~$S$. Mula himpunan-himpunan iki kawangun dadi
himpunan $\Setabs{x \in A}{\phi(x)}$ ing tataran~$S$
uga (miturut \stagesacc).

Beda karo Komprehensi Naif, prinsip iki nyingkiri Paradoks
Russell. Kita ora bisa langsung negesake anane himpunan
$\Setabs{x}{x \notin x}$. Nanging, yen \emph{diwenehi}
sawijining himpunan~$A$, kita bisa negesake anane himpunan
$R_A = \Setabs{x \in A}{x \notin x}$. Nanging kang bisa
dibuktekake mung $R_A \notin R_A$ lan $R_A \notin A$,
lan ora ana kang gawe kuwatir saka loro iki.

Nanging Pamisahan nduweni akibat langsung kang wigati:

\begin{thm}\ollabel{thm:NoUniversalSet}
Ora ana himpunan \emph{universal}, yaiku
$\Setabs{x}{x = x}$ ora ana.
\end{thm}

\begin{proof}
Kanggo bukti kanthi kontradiksi, upamane $V$ himpunan universal.
Miturut Pamisahan, mula $R =
\Setabs{x \in V}{x \notin x} = \Setabs{x}{x \notin x}$ ana,
lan iki lelawanan karo Paradoks Russell.
\end{proof}

Ora anane himpunan universal---malah sifat tanpa wates pungkasan
saka hierarki himpunan---minangka salah siji gagasan dhasar
ing konsepsi kumulatif-iteratif. Mula migunani yen kita
ndeleng yen kanthi intuitif asil iki uga bisa digayuh
liwat cara liyane. Himpunan universal mesthi dadi
!!a{element} saka awake dhewe. Nanging miturut konsepsi
kumulatif-iteratif, saben himpunan pisanan katon ing
tataran kang luwih pungkasan tinimbang saben !!{element}s-e.
Iki nuduhake yen \emph{ora ana} himpunan kang dadi anggota
awake dhewe. Amarga himpunan kaya mangkono kudu pisanan
katon ing tataran kang luwih pungkasan tinimbang tataran
pisanan nalika dheweke dhewe katon, lan iku mokal.
(Ing \olref[sth][spine][rank]{defnsetrank}, kita bakal
ndeleng carane nggawe panjelasan iki luwih trep nganggo
gagasan ``rank'' sawijining himpunan. Nanging kanggo iku
kita isih butuh sawetara aksioma liyane.)

Iki sawetara akibat liyane saka Pamisahan lan Ekstensionalitas.

\begin{prop}\ollabel{prop:emptyexists}
Yen ana himpunan apa wae, mula $\emptyset$ ana.
\end{prop}

\begin{proof}
Yen $A$ iku himpunan, mula miturut Pamisahan
$\emptyset = \Setabs{x \in A}{x \neq x}$ ana.
\end{proof}

\begin{prop}
$A \setminus B$ ana kanggo sembarang himpunan $A$ lan $B$.
\end{prop}

\begin{proof}
$A \setminus B = \Setabs{x \in A}{x \notin B}$ ana miturut Pamisahan.
\end{proof}

Uga bisa dibuktekake yen irisan saka (meh) sembarang kulawarga ana:

\begin{prop}\ollabel{prop:intersectionsexist}
Yen $A \neq \emptyset$, mula $\bigcap A = \Setabs{x}{(\forall y \in A)x \in y}$ ana.
\end{prop}

\begin{proof}
Upamane $A \neq \emptyset$, mula ana sawenehing $c \in A$.
Banjur $\bigcap A =
\Setabs{x}{(\forall y \in A)x \in y} = \Setabs{x \in c}{(\forall y \in
A)x \in y}$, kang ana miturut Pamisahan.
\end{proof}

Nanging gatekna syarat $A \neq \emptyset$; amarga $\bigcap
\emptyset$ bakal dadi himpunan universal kanthi vakum,
kang lelawanan karo \olref{thm:NoUniversalSet}.

\end{document}
